\documentclass[dvipsnames]{article}
\usepackage[english]{babel}
\usepackage{multirow}
\usepackage[letterpaper,top=2cm,bottom=2cm,left=3cm,right=3cm,marginparwidth=1.75cm]{geometry}
\usepackage{amsmath}
\usepackage{graphicx}
\usepackage{minted}
\usepackage{amsthm}
\usepackage{mathtools}
\usepackage[colorlinks=true, allcolors=blue]{hyperref}
\usepackage{fancyhdr}

\begin{document}
\section*{Week 1. Problem set (solutions by Karim Khabibrakhmanov)}

\thispagestyle{fancy}
\lhead{Karim Khabibrakhmanov (\texttt{k.khabibrakhmanov@innopolis.university}), group B24-DSAI-05}

\begin{enumerate}
  \item
    Compute asymptotic worst case time complexity of the following algorithm (see pseudocode conventions in \cite[\S 2.1]{Cormen2022}).
    You \textbf{must} use $\Theta$-notation.
    For justification, provide execution cost and frequency count for each line in the body of the \texttt{secret} procedure.
    Optionally, you may provide the details for the computation of the running time $T(n)$ for worst case scenario.
    Proof for the asymptotic bound is not required for this exercise.

    \begin{minipage}{0.55\textwidth}
      \begin{minted}[linenos,escapeinside=||]{sql}
/* A is a 0-indexed array,
  * n is the number of items in A */
|\color{blue}secret|(A, n):
  |\color{red}for| i = 1 to n
    s := 1
    j := i - 1
    |\color{red}while| j < n |\color{teal}and| A[j] |$\geq$| A[i-1]
      j := j + s
      s := s + 2
    |\color{teal}exchange| A[i-1] |\color{teal}with| A[min(j, n - 1)]
      \end{minted}
    \end{minipage}%
    \begin{minipage}{0.4\textwidth}
      \small
      \begin{tabular}{ c | c | c }
      Line & Execution Cost & Frequency Count \\ 
      \hline
      1 & --- & --- \\  
      2 & --- & --- \\
      3 & --- & --- \\
      4 &     &  \\
      5 &     &  \\
      6 &     &  \\
      7 &     &  \\
      8 &     &  \\
      9 &     &  \\
      10 &    &  \\
      \end{tabular}
    \end{minipage}

\color{blue}
\textbf{Answer:} \emph{$\Theta(n^2)$}. \\
\textbf{Justification: } \emph{cost and frequency count is presented in the following table:} \\
    \begin{minipage}{0.4\textwidth}
      \small
      \begin{tabular}{ c | c | c }
      Line & Execution Cost & Frequency Count \\ 
      \hline
      1 & 0 & 1 \\  
      2 & 0 & 1 \\
      3 & $c_0$ & 1 \\
      4 &  $c_1$   & n+1 \\
      5 &  $c_2$   & n \\
      6 &  $c_3$   & n \\
      7 &  $c_4$   & \[ \sum_{i=1}^{n} t_i \] \\
      8 &  $c_5$  & \[ \sum_{i=1}^{n} (t_i-1) \] \\
      9 &   $c_6$  & \[ \sum_{i=1}^{n} (t_i-1) \] \\
      10 &  $c_7$  & n \\
      \end{tabular}
    \end{minipage}

    \emph{(optional part) From this table, we conclude}
$$T(n) = c_0 \cdot 1 + c_1 \cdot (n+1) + c_2 \cdot n + c_3 \cdot n + c_4 \cdot \sum_{i=1}^{n} t_i + c_5 \cdot \sum_{i=1}^{n} (t_i-1) + c_6 \cdot \sum_{i=1}^{n} (t_i-1) + c_7 \cdot n = \Theta(n^2)$$

\color{black}
  \item Indicate, for each pair of expressions (A, B) in the table below whether $A = O(B)$, $A = o(B)$, $A = \Omega(B)$, $A = \omega(B)$, or $A = \Theta(B)$. Write your answer in the form of the table with \emph{YES} or \emph{NO} written in each cell:

    \begin{center}
    \bgroup
    \def\arraystretch{1.5}
    \begin{tabular}{ |c|c||c|c|c|c|c| }
     \hline
     $A$ & $B$ & $A = O(B)$ & $A = o(B)$ & $A = \Omega(B)$ & $A = \omega(B)$ & $A = \Theta(B)$ \\
     \hline
     \hline
     $1 + \sin(n)$ & $\left(1 + \frac{1}{n}\right)^{\frac{1}{n}}$ & \textcolor{blue}{\emph{YES}} & \textcolor{blue}{\emph{NO}} & \textcolor{blue}{\emph{NO}} & \textcolor{blue}{\emph{NO}} & \textcolor{blue}{\emph{NO}} \\
     \hline
     $\sqrt[5]{n}$ & $\log_3^2 n$ & \textcolor{blue}{\emph{NO}} & \textcolor{blue}{\emph{NO}} & \textcolor{blue}{\emph{YES}}  & \textcolor{blue}{\emph{YES}}  & \textcolor{blue}{\emph{NO}} \\
     \hline
     $n^2$ & $\log_2^n n$ & \textcolor{blue}{\emph{YES}} & \textcolor{blue}{\emph{YES}} & \textcolor{blue}{\emph{NO}} & \textcolor{blue}{\emph{NO}} & \textcolor{blue}{\emph{NO}} \\
     \hline
     $(1 + \frac{1}{10^{100}})^n$ & $n^{(10^{100})}$ & \textcolor{blue}{\emph{NO}} & \textcolor{blue}{\emph{NO}} & \textcolor{blue}{\emph{YES}} & \textcolor{blue}{\emph{YES}} & \textcolor{blue}{\emph{NO}} \\
     \hline
     $\log_2 \left( (n + 1)! \right)$ & $n \log_2 n$ & \textcolor{blue}{\emph{YES}} & \textcolor{blue}{\emph{NO}} & \textcolor{blue}{\emph{YES}} & \textcolor{blue}{\emph{NO}} & \textcolor{blue}{\emph{YES}} \\
     \hline
    \end{tabular}
    \egroup
    \end{center}

  \clearpage
  \item Let $f$ and $g$ be functions from positive integers to positive reals.
    Assume $f(n) > n$ for $n > 2025$.
    Using the formal definition of asymptotic notation~\cite[\S 3.2]{Cormen2022}, prove \emph{formally} that
    \[
      \min(f(n) - \sqrt{n}, g(n) + n) = O(f(n) + g(n))
    \]
  \color{blue}
  \begin{proof}
    \emph{By definition of big-$O$ notation, we need to show that there exist constant c and $n_0$ such that for all n $ \geq $ $n_0$}
    \[
      \min(f(n) - \sqrt{n}, g(n) + n) \leq \ c * (f(n) + g(n))
    \]

    \emph{Let c = $1$ and $n_0$ = $2025$}

    \emph{By assumption, we have $f(n) > n$ for $n > 2025$}

    \emph{Consider the following cases:}\\
    \hspace*{3mm}
    1) if $(f(n) - \sqrt{n}) \geq (g(n) + n)$:\\
    \hspace*{10mm}
    $g(n) + n \leq\ 1 * (f(n) + g(n))$\\
    \hspace*{10mm}
    $g(n) + n \leq\ 1 * f(n) + 1 * g(n)$\\
    \hspace*{8mm}
    Here we subtract $g(n)$ from both sides.\\
    \hspace*{10mm}
    $n \leq\ f(n)$\\
    \hspace*{8mm}
    Since we have established that $ f(n) > n $ for $ n > 2025 $, the inequality holds true.\\ 
    \hspace*{8mm}
    Thus, the condition $ n \leq f(n) $ is satisfied for $ n > 2025 $.\\

    
    \hspace*{3mm}
    2) if $(f(n) - \sqrt{n}) \leq (g(n) + n)$:\\
    \hspace*{10mm}
    $f(n) - \sqrt{n} \leq\ 1 * (f(n) + g(n))$\\
    \hspace*{10mm}
    $f(n) - \sqrt{n} \leq\ 1 * f(n) + 1 * g(n)$\\
    \hspace*{8mm}
    Here we subtract $f(n)$ from both sides.\\
    \hspace*{10mm}
    $- \sqrt{n} \leq\ g(n)$\\
    \hspace*{8mm}
    Since $ g(n) $ is a positive function and $ n > 2025 $, it follows that $ \sqrt{n} $ is also positive. Therefore, \\
    \hspace*{8mm}
    $ -\sqrt{n} $ is negative. This implies that the inequality $ -\sqrt{n} \leq g(n) $ holds true, as a negative\\
    \hspace*{8mm}
    value is always less than or equal to a positive value. Thus, the condition is satisfied.\\

    
    \hspace*{3mm}    
    3) Conclusion:\\
    \hspace*{8mm}
    We have examined all possible cases for Big-O, and both have been found to be true.\\
    \hspace*{8mm}
    Therefore, the statement $ \min(f(n) - \sqrt{n}, g(n) + n) = O(f(n) + g(n)) $ is also true.\\

    \emph {QED.}\\


  \end{proof}
  \color{black}

\end{enumerate}

\begin{thebibliography}{TheLongestRefName}

\bibitem[CLRS]{Cormen2022}
Cormen, T.H., Leiserson, C.E., Rivest, R.L. and Stein, C., 2022. \emph{Introduction to algorithms, Fourth Edition}. MIT press.

%\bibitem[Goldrich]{Goldrich2014}
%M. T. Goodrich, R. Tamassia, and M. H. Goldwasser. \emph{Data Structures and Algorithms in Java.} WILEY 2014.
\end{thebibliography}

\end{document}
